% MATA-LL Developer's Manual
% Compile with: lualatex devmanual.tex
\documentclass[11pt,a4paper]{article}

\usepackage{fontspec}
\usepackage{unicode-math}
\setmainfont{Latin Modern Roman}
\setmonofont{Latin Modern Mono}[Scale=0.88]
\setmathfont{Latin Modern Math}

\usepackage{geometry}
\geometry{margin=2.5cm}

\usepackage{hyperref}
\hypersetup{colorlinks=true, linkcolor=blue!60!black, urlcolor=blue!60!black}

\usepackage{listings}
\usepackage{xcolor}
\usepackage{booktabs}
\usepackage{enumitem}
\usepackage{titlesec}

\titleformat{\section}{\Large\bfseries}{\thesection}{1em}{}
\titleformat{\subsection}{\large\bfseries}{\thesubsection}{1em}{}

\definecolor{codebg}{HTML}{F5F5F5}
\definecolor{keyword}{HTML}{7B3294}
\definecolor{strcolor}{HTML}{2E7D32}
\definecolor{comment}{HTML}{757575}

\lstdefinelanguage{MLL}{
  keywords={data,where,let,in,case,of,if,then,else,do,import,module,
            class,instance,deriving,newtype,type,export,infixl,infixr,infix,forall},
  sensitive=true,
  morecomment=[l]{--},
  morecomment=[n]{\{-}{-\}},
  morestring=[b]",
}

\lstset{
  basicstyle=\ttfamily\small,
  backgroundcolor=\color{codebg},
  frame=single,
  framerule=0pt,
  rulecolor=\color{codebg},
  keywordstyle=\color{keyword}\bfseries,
  stringstyle=\color{strcolor},
  commentstyle=\color{comment}\itshape,
  breaklines=true,
  showstringspaces=false,
  tabsize=2,
  xleftmargin=4pt,
  xrightmargin=4pt,
  aboveskip=8pt,
  belowskip=8pt,
}

\lstdefinelanguage{Rust}{
  keywords={fn,pub,struct,enum,impl,let,mut,use,mod,match,if,else,
            return,self,Self,trait,where,type,const,for,in,loop,while},
  sensitive=true,
  morecomment=[l]{//},
  morecomment=[n]{/*}{*/},
  morestring=[b]",
}

\title{\textbf{MATA-LL Developer's Manual}\\[6pt]
  \large Modest Attempt at Typesystem Augmenting the Lua Language}
\author{MATA-LL Project}
\date{June 2026}

\begin{document}
\maketitle
\tableofcontents
\clearpage

%----------------------------------------------------------------------
\section{Introduction}

MATA-LL is a Haskell-inspired language that compiles to standalone Lua files with
no runtime dependencies.  The compiler is written in Rust (\textasciitilde15\,000
lines) and targets both Lua~5.4 and LuaJIT.

\begin{itemize}[nosep]
  \item File extension: \texttt{.mll}
  \item Output: a single self-contained \texttt{.lua} file
  \item Non-strict evaluation with memoizing thunks
  \item Full Hindley--Milner type inference with typeclasses, GADTs, rank-2 types
  \item Zero external Rust dependencies in the core compiler library
\end{itemize}

%----------------------------------------------------------------------
\section{Repository Layout}

\begin{tabular}{@{}ll@{}}
\toprule
\textbf{Path} & \textbf{Purpose} \\
\midrule
\texttt{mllc/}       & Compiler library (core) \\
\texttt{mll/}        & CLI binary \\
\texttt{mll-repl/}   & Interactive REPL \\
\texttt{mll-tests/}  & Test harness (246 tests) \\
\texttt{mllc-wasm/}  & WASM binding for browser playground \\
\texttt{lib/}        & Standard library (\texttt{.mll} files) \\
\texttt{examples/}   & Example programs \\
\texttt{doc/}        & This manual \\
\bottomrule
\end{tabular}

%----------------------------------------------------------------------
\section{Building}

\begin{lstlisting}[language=bash]
cargo build --release        # all crates
cargo test                   # run 246 tests (needs 32 MB stack)
cargo run -- file.mll --run  # compile and execute
\end{lstlisting}

The CLI binary (\texttt{mll}) embeds \texttt{mlua} with vendored Lua~5.4 for
\texttt{-{}-run} mode.  The REPL uses LuaJIT via \texttt{mlua}.

%----------------------------------------------------------------------
\section{Compilation Pipeline}

The compiler transforms source through seven distinct phases:

\begin{enumerate}[nosep]
  \item \textbf{Lexing} --- layout-sensitive tokenization
  \item \textbf{Parsing} --- recursive descent with operator-precedence climbing
  \item \textbf{Import resolution} --- module loading and cycle detection
  \item \textbf{Desugaring} --- do-notation to monadic bind chains
  \item \textbf{Type checking} --- HM unification + bidirectional checking
  \item \textbf{Monomorphization} --- specialization + typeclass dispatch
  \item \textbf{Demand analysis} --- strictness inference
  \item \textbf{Code generation} --- Lua source emission
\end{enumerate}

The public API is a single function:
\begin{lstlisting}[language=Rust]
pub fn compile(
    source: &str,
    source_dir: &Path,
    lib_paths: &[&Path],
) -> Result<CompileResult, CompileError>
\end{lstlisting}

Returning \texttt{CompileResult \{ lua\_code, has\_main, exports, warnings \}}.
The \texttt{warnings} are non-fatal diagnostics the frontend should surface;
currently: the module has neither \texttt{main} nor any \texttt{export}, so
dead-code elimination (rooted at exactly those two) removed every definition
and the emitted Lua contains nothing to run or call.  A module-header export
list (\texttt{module M (foo) where}) does not prevent this --- it scopes
\texttt{.mll}-level import visibility only and never populates the Lua
export table.

%----------------------------------------------------------------------
\subsection{Lexer (\texttt{mllc/src/lexer.rs})}

The lexer is layout-sensitive.  It emits \texttt{Indent(n)} at the start of each
indented line and \texttt{Newline} tokens between logical lines.  Block comments
\texttt{\{- \ldots\ -\}} nest correctly.  Operators are maximal sequences of
symbol characters.

%----------------------------------------------------------------------
\subsection{Parser (\texttt{mllc/src/parser.rs})}

A hand-written recursive-descent parser with operator-precedence climbing for
infix expressions.  Fixity is tracked in a \texttt{HashMap<String, (Assoc, u8)>}
and respects user \texttt{infixl}/\texttt{infixr}/\texttt{infix} declarations.

Layout handling: continuation lines indented past the definition's start column
are merged into the same declaration.

Declaration forms: \texttt{TypeSig}, \texttt{FunDef}, \texttt{DataDef},
\texttt{NewtypeDef}, \texttt{ClassDecl}, \texttt{InstanceDecl},
\texttt{TypeFamily}, \texttt{TypeAlias}, \texttt{FixityDecl}, \texttt{Import},
\texttt{ExportSig}.

%----------------------------------------------------------------------
\subsection{Import Resolution (\texttt{mllc/src/modules.rs})}

Module paths map dots to directories: \texttt{Data.Map} resolves to
\texttt{Data/Map.mll}.  Search order: source directory, then each path supplied
via \texttt{-L}.  Loaded modules are cached and cycles are detected.

The Prelude is embedded in the binary via \texttt{include\_str!} and
prepended automatically before desugaring.

%----------------------------------------------------------------------
\subsection{Desugarer (\texttt{mllc/src/desugar.rs})}

Eliminates do-notation:
\begin{lstlisting}[language=MLL]
do { x <- action; rest }   -- becomes: action >>= \x -> rest
do { action; rest }         -- becomes: action >>= \_ -> rest
do { let x = e; rest }     -- becomes: let x = e in rest
\end{lstlisting}

%----------------------------------------------------------------------
\subsection{Type Checker (\texttt{mllc/src/typechecker.rs})}

Robinson unification with bidirectional checking.  Top-level definitions require
explicit type signatures; local bindings are inferred.  Errors accumulate rather
than aborting on the first failure.

Key capabilities:
\begin{itemize}[nosep]
  \item Kind system: \texttt{Type}, \texttt{Symbol}, \texttt{Arrow(k1, k2)}
  \item Typeclasses with superclass constraints and default methods
  \item GADTs with local type equalities introduced by pattern matching
  \item Rank-2 types (\texttt{forall s.}) for scoped resources
  \item Existential types in data constructors
  \item Type families (intrinsic and user-defined, closed)
  \item DataKinds (promoted constructors as type-level tags)
  \item Deriving: \texttt{Show}, \texttt{Eq}, \texttt{Ord}, \texttt{Enum},
        \texttt{Bounded}, \texttt{Functor}
\end{itemize}

%----------------------------------------------------------------------
\subsection{Monomorphizer (\texttt{mllc/src/mono.rs})}

Collects \texttt{(function\_name, concrete\_type)} pairs and generates
specialized copies with mangled names.  Typeclass method calls are resolved to
concrete instance functions.

When specialization count for a function exceeds 16, the monomorphizer falls back
to dictionary-passing: the function receives a Lua table of method
implementations as a parameter.

%----------------------------------------------------------------------
\subsection{Demand Analysis (\texttt{mllc/src/demand.rs})}

A fixed-point analysis with cross-function propagation.  Determines which
parameters are forced on every code path.  Strict parameters skip thunk
allocation at call sites and \texttt{\_\_force} at use sites.  FFI functions are
seeded as strict in all parameters.

%----------------------------------------------------------------------
\subsection{Code Generator (\texttt{mllc/src/codegen.rs})}

Emits Lua source.  Key design decisions:
\begin{itemize}[nosep]
  \item Forward declarations packed into an \texttt{\_\_mll\_fn} table (avoids
        Lua's 200-local limit)
  \item Local overflow uses a \texttt{\_v = \{\}} table when count exceeds
        \textasciitilde180
  \item Cheapness analysis determines whether expressions become thunks or
        are evaluated eagerly
  \item Small pure single-clause non-recursive functions are inlined
  \item Bind chains (do-blocks) are flattened to sequential locals
  \item Record field accessors become \texttt{\_r[N]} inline access
\end{itemize}

%----------------------------------------------------------------------
\section{Runtime Representation}

\begin{tabular}{@{}ll@{}}
\toprule
\textbf{MLL construct} & \textbf{Lua representation} \\
\midrule
Multi-constructor ADT \texttt{Con x y}  & \texttt{\{tag, x, y\}} \\
Single-constructor, single-field        & Unwrapped (no table) \\
Newtype                                 & Identity (zero-cost) \\
Pure enum (\texttt{Red | Green | Blue}) & Integers 1, 2, 3 \\
\texttt{Nothing}                        & \texttt{nil} \\
\texttt{Just x}                         & \texttt{x} directly \\
Tuple \texttt{(a, b, c)}                & \texttt{\{a, b, c\}} (no tag) \\
List \texttt{x : xs}                    & \texttt{\{x, xs\}} cons cell; \texttt{[]} = \texttt{nil} \\
Thunk                                   & \texttt{\{fn, false\}} with \texttt{\_\_thunk\_mt} \\
\texttt{HashMap}                        & Native Lua table \\
\bottomrule
\end{tabular}

\medskip

Thunk forcing is handled by a runtime preamble embedded in every generated file:
\begin{lstlisting}[language={}]
local function __thunk(f)
  return setmetatable({f, false}, __thunk_mt)
end
local function __force(x)
  if getmetatable(x) == __thunk_mt then
    if x[2] then return x[1] end
    local val = x[1](); x[1] = val; x[2] = true
    return val
  end
  return x
end
\end{lstlisting}

%----------------------------------------------------------------------
\section{Type System}

\subsection{Core System}

Hindley--Milner with let-generalization.  Unification variables are \texttt{u32}
IDs; skolem/rigid variables use \texttt{u32::MAX} range.  Substitutions are
applied iteratively (not recursively) to avoid stack overflow on deeply nested
bind chains.

\subsection{Typeclasses}

Single-parameter typeclasses with superclass constraints.  Built-in classes:
\texttt{Monad}, \texttt{Show}, \texttt{Eq}, \texttt{Ord}, \texttt{Enum},
\texttt{Bounded}, \texttt{Read}, \texttt{Functor}, \texttt{Applicative}.

Instance resolution key: \texttt{(class\_name, type\_name)} mapping to method
function names.  Orphan instances are rejected.

\subsection{GADTs}

Constructors carry full return types.  Pattern matching introduces local type
equalities via unification.  Exhaustiveness checking accounts for GADT
refinement.

\subsection{Rank-2 Types}

\texttt{forall s.} quantification specifically for:
\begin{itemize}[nosep]
  \item \texttt{LuaFunction s} --- scope safety for Lua callbacks
  \item \texttt{runST :: (forall s. ST s a) -> a} --- scoped mutation
\end{itemize}

\subsection{Type Families}

Intrinsic type families for FFI:
\begin{itemize}[nosep]
  \item \texttt{LuaPure "name" a} --- pure foreign function
  \item \texttt{LuaIO "name" a} --- effectful foreign function
  \item \texttt{LuaIterator "name" a} --- Lua iterator protocol
  \item \texttt{LuaTry "name" a} --- pcall-wrapped call
\end{itemize}

User-defined closed type families (equation-based reduction) are also supported.

\subsection{DataKinds}

Promoted constructors serve as type-level tags:
\begin{lstlisting}[language=MLL]
data Safety = Safe | Unsafe
data Handle (s :: Safety) = MkHandle String

lock :: Handle 'Unsafe -> Handle 'Safe
\end{lstlisting}

%----------------------------------------------------------------------
\section{FFI}

Foreign functions are declared via type families:
\begin{lstlisting}[language=MLL]
sin :: LuaPure "math.sin" (Number -> Number)
randomInt :: LuaIO "math.random" (Integer -> Integer -> Integer)
\end{lstlisting}

Method-call syntax uses a colon prefix:
\begin{lstlisting}[language=MLL]
fileWrite :: LuaIO ":write" (Handle -> String -> IO ())
\end{lstlisting}

A top-level \texttt{Maybe} parameter becomes an optional Lua argument.  Outgoing
call arguments are marshalled by \texttt{\_\_mll\_arg\_marshal}, the structural
dual of the result decoder \texttt{\_\_mll\_ffi\_decode}: it descends into lists,
tuples, \texttt{LuaDict} records, \texttt{HashMap} values, and \texttt{Maybe}
(unwrapping \texttt{Just}/\texttt{Nothing}), rebuilding fresh Lua values --- so
any type the decoder converts on the way in is converted on the way out --- while
leaving opaque values (type variables, plain ADTs, functions) raw.  The
\texttt{export} keyword wraps return values via \texttt{\_\_mll\_to\_lua}
(deep-forces thunks, converts cons lists to Lua arrays) and incoming callbacks
via \texttt{\_\_mll\_wrap\_callback}.

%----------------------------------------------------------------------
\section{Standard Library}

\begin{tabular}{@{}ll@{}}
\toprule
\textbf{Module} & \textbf{File} \\
\midrule
Prelude (auto-imported) & \texttt{lib/Prelude.mll} \\
\texttt{ByteString}     & \texttt{lib/ByteString.mll} \\
\texttt{LIO}            & \texttt{lib/LIO.mll} \\
\texttt{LMath}          & \texttt{lib/LMath.mll} \\
\texttt{LOS}            & \texttt{lib/LOS.mll} \\
\texttt{LString}        & \texttt{lib/LString.mll} \\
\texttt{LBit}           & \texttt{lib/LBit.mll} \\
\texttt{Regex}          & \texttt{lib/Regex.mll} \\
\texttt{JSON}           & \texttt{lib/JSON.mll} \\
\texttt{Data.List}      & \texttt{lib/Data/List.mll} \\
\texttt{Data.Maybe}     & \texttt{lib/Data/Maybe.mll} \\
\texttt{Data.Map}       & \texttt{lib/Data/Map.mll} \\
\texttt{Control.Monad}  & \texttt{lib/Control/Monad.mll} \\
\bottomrule
\end{tabular}

%----------------------------------------------------------------------
\section{Deviations from Haskell}

Key intentional differences:
\begin{itemize}[nosep]
  \item \texttt{String} is a native Lua string, not \texttt{[Char]}.
        No \texttt{Char} type exists.
  \item No \texttt{Num} typeclass --- arithmetic uses built-in operators on
        \texttt{Integer} and \texttt{Number} directly.
  \item No multi-parameter typeclasses.
  \item No lazy I/O --- file operations read eagerly.
  \item No \texttt{IORef}, \texttt{MVar}, \texttt{STRef} ---
        \texttt{STArray} covers scoped mutation.
  \item String concatenation uses \texttt{<>} (Semigroup), not \texttt{++}
        (which is list append).
\end{itemize}

%----------------------------------------------------------------------
\section{Testing}

Tests live in \texttt{mll-tests/tests/cases/} (171 case files) and
\texttt{mll-tests/tests/ghc/} (GHC compatibility tests).

\begin{lstlisting}[language=Rust]
mll_test!(basics, "basics.mll");            // no lib path
mll_lib_test!(json_parse, "json.mll");      // with ../lib
\end{lstlisting}

Each test file uses an \texttt{assert :: Bool -> String -> IO ()} that prints
\texttt{"."} on success or calls \texttt{error} on failure.

Run all tests:
\begin{lstlisting}[language=bash]
cargo test -p mll-tests
\end{lstlisting}

Tests run on 32\,MB stack threads (deeply nested ASTs from large list literals).

%----------------------------------------------------------------------
\section{Adding a New Feature: Checklist}

\begin{enumerate}[nosep]
  \item \textbf{Lexer}: add new tokens if needed (\texttt{lexer.rs})
  \item \textbf{AST}: add node variants (\texttt{ast.rs})
  \item \textbf{Parser}: handle new syntax (\texttt{parser.rs})
  \item \textbf{Desugarer}: add rewrite if the feature is syntactic sugar
  \item \textbf{Type checker}: extend inference/checking (\texttt{typechecker.rs})
  \item \textbf{TIR}: add typed IR node if needed (\texttt{tir.rs})
  \item \textbf{Monomorphizer}: handle new TIR node (\texttt{mono.rs})
  \item \textbf{Demand analysis}: propagate strictness if relevant
  \item \textbf{Codegen}: emit Lua (\texttt{codegen.rs})
  \item \textbf{Tests}: add case file + \texttt{mll\_test!} entry
\end{enumerate}

%----------------------------------------------------------------------
\section{Source Map}

\begin{tabular}{@{}lrl@{}}
\toprule
\textbf{File} & \textbf{Lines} & \textbf{Purpose} \\
\midrule
\texttt{mllc/src/lib.rs}          & 126  & Public API, pipeline orchestration \\
\texttt{mllc/src/lexer.rs}        & 443  & Layout-sensitive tokenizer \\
\texttt{mllc/src/ast.rs}          & 352  & AST type definitions \\
\texttt{mllc/src/parser.rs}       & 2481 & Recursive-descent parser \\
\texttt{mllc/src/modules.rs}      & 237  & Import resolution \\
\texttt{mllc/src/desugar.rs}      & 195  & Do-notation desugaring \\
\texttt{mllc/src/typechecker.rs}  & 3944 & Type inference and checking \\
\texttt{mllc/src/types.rs}        & 496  & Internal type representation \\
\texttt{mllc/src/tir.rs}          & 356  & Typed intermediate representation \\
\texttt{mllc/src/mono.rs}         & 1534 & Monomorphization \\
\texttt{mllc/src/demand.rs}       & 394  & Demand/strictness analysis \\
\texttt{mllc/src/codegen.rs}      & 3120 & Lua code generation \\
\bottomrule
\end{tabular}

%----------------------------------------------------------------------
\section{CLI Reference}

\begin{lstlisting}[language=bash]
mll file.mll              # compile, print Lua to stdout
mll file.mll --run        # compile and execute
mll file.mll --emit-lua   # write file.lua
mll file.mll -L ./lib     # add library search path
\end{lstlisting}

REPL commands:
\begin{tabular}{@{}ll@{}}
\toprule
\textbf{Command} & \textbf{Effect} \\
\midrule
\texttt{/clear}  & Reset all accumulated state \\
\texttt{/decls}  & Show current declarations \\
\texttt{/lua EXPR} & Evaluate raw Lua expression \\
\texttt{/source} & Show generated Lua source \\
\texttt{/drop}   & Remove last declaration \\
\texttt{/help}   & Show help \\
\texttt{/quit}   & Exit \\
\bottomrule
\end{tabular}

\medskip
Multi-line input: end a line with \texttt{\textbackslash} to continue on the next line.

\end{document}
